
\dangbt{chưa phân loại}
\dangbt{Vận dụng đồ thị độ dịch chuyển -- thời gian của một vật}
\begin{ex}
Hình bên mô tả đồ thị tọa độ - thời gian của hai chiếc xe trong cùng một khoảng thời gian. Xét các phát biểu sau:
\choiceTF
{Hai xe chuyển động thẳng đều, ngược chiều nhau}
{Với cùng một khoảng thời gian, xe 1 đi được quãng đường lớn hơn xe 2}
{\True Tốc độ tức thời của xe 2 lớn hơn xe 1}
{Xe 1 có vận tốc tức thời lớn hơn xe 2}
\loigiai{
a sai vì chưa đủ dữ kiện xác định cả hai đều ngược chiều nhau; b sai vì độ dốc đồ thị xe 2 lớn hơn nên trong cùng khoảng thời gian xe 2 đi được quãng đường lớn hơn xe 1; c đúng vì độ lớn hệ số góc của đồ thị xe 2 lớn hơn xe 1 nên tốc độ tức thời của xe 2 lớn hơn; d sai vì cần xét cả dấu của hệ số góc trên đồ thị tọa độ - thời gian.
}
\end{ex}

\dangbt{Vận dụng đồ thị độ dịch chuyển -- thời gian của một vật}
\begin{ex}
Một vật chuyển động thẳng có đồ thị độ dịch chuyển – thời gian $(d-t)$. Xét các phát biểu sau:
\choiceTF
{\True Tại $A$, vật chuyển động thẳng đều, cùng chiều dương}
{Tốc độ tức thời ở vị trí $A$ là $1\text{ m/s}$}
{\True Tại vị trí $B$, vật không chuyển động}
{Tại vị trí $C$, vật chuyển động thẳng chậm dần đều với tốc độ tức thời là $2\text{ m/s}$}
\loigiai{
a đúng vì đoạn đồ thị tại $A$ là đường thẳng dốc lên theo chiều dương; b sai vì tốc độ tức thời tính theo độ dốc khác giá trị $1\text{ m/s}$; c đúng vì tại đoạn chứa điểm $B$ đồ thị song song với trục thời gian nên vật đứng yên; d sai vì đoạn thẳng thể hiện chuyển động thẳng đều chứ không phải chậm dần đều.
}
\end{ex}

\dangbt{Vận dụng đồ thị độ dịch chuyển -- thời gian của một vật}
\begin{ex}
Đồ thị tọa độ – thời gian có dạng một nhánh parabol $x(t) = k t^2$ ($k > 0$). Chuyển động của vật là:
\choice
{Chuyển động thẳng đều}
{Chuyển động chậm dần đều}
{\True Chuyển động nhanh dần đều}
{Vận tốc không đổi}
\loigiai{
Đồ thị $x(t)$ có dạng parabol bậc hai theo thời gian ứng với chuyển động biến đổi đều có vận tốc tăng dần theo thời gian, tức là chuyển động nhanh dần đều.
}
\end{ex}

\dangbt{Vận dụng đồ thị độ dịch chuyển -- thời gian của một vật}
\begin{ex}
Đồ thị tọa độ – thời gian $x(t)$ của một vật là một đường thẳng dốc lên theo thời gian. Phát biểu nào sau đây đúng về chuyển động của vật?
\choice
{Chuyển động đều theo chiều âm}
{Chuyển động chậm dần đều}
{\True Chuyển động thẳng đều}
{Chuyển động nhanh dần đều}
\loigiai{
Đồ thị $x(t)$ là đường thẳng có hệ số góc không đổi nên vận tốc của vật là hằng số, tức là vật chuyển động thẳng đều.
}
\end{ex}

\dangbt{Vận dụng đồ thị độ dịch chuyển -- thời gian của một vật}
\begin{ex}
Cho đồ thị dịch chuyển – thời gian của một vật gồm: đoạn từ $0$ đến $t_1$ là đường thẳng dốc lên, đoạn từ $t_1$ đến $t_2$ nằm ngang song song với trục thời gian, đoạn từ $t_2$ đến $t_3$ là đường thẳng dốc xuống. Trong những khoảng thời gian nào, vật chuyển động thẳng đều?
\choice
{Trong khoảng thời gian từ 0 đến $t_1$ và từ $t_1$ đến $t_2$}
{Trong khoảng thời gian từ $t_1$ đến $t_2$}
{Trong khoảng thời gian từ 0 đến $t_3$}
{\True Trong khoảng thời gian từ 0 đến $t_1$ và từ $t_2$ đến $t_3$}
\loigiai{
Trong khoảng thời gian từ $t_1$ đến $t_2$, đồ thị nằm ngang nên vật đứng yên. Các khoảng từ $0$ đến $t_1$ và từ $t_2$ đến $t_3$ có đồ thị là đường thẳng nghiêng nên vật chuyển động thẳng đều.
}
\end{ex}
\dangbt{Vận dụng đồ thị độ dịch chuyển -- thời gian của một vật}
\begin{ex}
  Một vật chuyển động thẳng có đồ thị độ dịch chuyển – thời gian như hình vẽ.
  
  (Hình vẽ: Đồ thị $x-t$. Trục hoành $t$ (s), trục tung $x$ (m).
  Điểm (0,0) đến (2, 4).
  Điểm (2, 4) đến (4, 4).
  Điểm (4, 4) đến (6, 0).
  Điểm (6, 0) đến (8, -4).
  Điểm (8, -4) đến (10, -4).)
  
  Tính vận tốc của vật trong khoảng thời gian từ $t = 0$ s đến $t = 2$ s (theo m/s).
  \shortans{2}
  \loigiai{
  Trong khoảng thời gian từ $t = 0$ s đến $t = 2$ s, đồ thị là một đoạn thẳng.
  Vận tốc của vật trong khoảng thời gian này là độ dốc của đoạn thẳng đó.
  $v = \frac{\Delta x}{\Delta t} = \frac{x_2 - x_0}{t_2 - t_0} = \frac{4 - 0}{2 - 0} = \frac{4}{2} = 2$ m/s.
  }
\end{ex}

\dangbt{Vận dụng đồ thị độ dịch chuyển -- thời gian của một vật}
\begin{ex}
Một vật chuyển động có đồ thị độ dịch chuyển – thời gian như hình vẽ.
Trong khoảng thời gian từ $t = 0$ đến $t = 2$ s, vật chuyển động thẳng đều.
Trong khoảng thời gian từ $t = 2$ s đến $t = 4$ s, vật đứng yên.
Trong khoảng thời gian từ $t = 4$ s đến $t = 6$ s, vật chuyển động thẳng đều.
Xác định vận tốc của vật trong khoảng thời gian từ $t = 4$ s đến $t = 6$ s (theo m/s).
\shortans{2}
\loigiai{
Trong khoảng thời gian từ $t = 4$ s đến $t = 6$ s:
Độ dịch chuyển tại $t = 4$ s là $x_1 = 4$ m.
Độ dịch chuyển tại $t = 6$ s là $x_2 = 8$ m.
Vận tốc của vật là $v = \frac{x_2 - x_1}{t_2 - t_1} = \frac{8 - 4}{6 - 4} = \frac{4}{2} = 2$ m/s.
}
\end{ex}

\dangbt{Vận dụng đồ thị độ dịch chuyển -- thời gian của một vật}
\begin{ex}
Một vật chuyển động thẳng có đồ thị độ dịch chuyển – thời gian như hình vẽ.
\begin{center}
\includegraphics[width=0.7\textwidth]{example_graph.png}
\end{center}
a) Mô tả chuyển động của vật trong từng giai đoạn (từ 0 đến 2s, từ 2s đến 4s, từ 4s đến 6s).
b) Tính vận tốc của vật trong từng giai đoạn.
c) Tính tổng độ dịch chuyển của vật sau 6s.
d) Vẽ đồ thị vận tốc – thời gian của vật.

\loigiai{
a) Mô tả chuyển động:
- Giai đoạn từ 0 đến 2s: Đồ thị là đoạn thẳng đi lên từ gốc tọa độ, chứng tỏ vật chuyển động thẳng đều theo chiều dương, bắt đầu từ vị trí 0m.
- Giai đoạn từ 2s đến 4s: Đồ thị là đoạn thẳng nằm ngang, chứng tỏ độ dịch chuyển không đổi, vật đứng yên tại vị trí 4m.
- Giai đoạn từ 4s đến 6s: Đồ thị là đoạn thẳng đi xuống, chứng tỏ vật chuyển động thẳng đều theo chiều âm, từ vị trí 4m về vị trí 0m.

b) Tính vận tốc:
- Giai đoạn từ 0 đến 2s: $v_1 = \frac{\Delta d_1}{\Delta t_1} = \frac{4 - 0}{2 - 0} = 2 \text{ m/s}$.
- Giai đoạn từ 2s đến 4s: $v_2 = \frac{\Delta d_2}{\Delta t_2} = \frac{4 - 4}{4 - 2} = 0 \text{ m/s}$.
- Giai đoạn từ 4s đến 6s: $v_3 = \frac{\Delta d_3}{\Delta t_3} = \frac{0 - 4}{6 - 4} = -2 \text{ m/s}$.

c) Tổng độ dịch chuyển sau 6s:
Độ dịch chuyển cuối cùng của vật tại $t = 6s$ là $d = 0 \text{ m}$.
Vậy tổng độ dịch chuyển của vật sau 6s là $0 \text{ m}$.

d) Đồ thị vận tốc – thời gian:
\begin{center}
\includegraphics[width=0.7\textwidth]{velocity_time_graph.png}
\end{center}
}
\end{ex}

\dangbt{Vận dụng đồ thị độ dịch chuyển -- thời gian của một vật}
\begin{ex}
Một xe máy chuyển động thẳng có đồ thị độ dịch chuyển – thời gian như hình vẽ.
\begin{center}
\includegraphics[width=0.7\textwidth]{graph_displacement_time_2.png}
\end{center}
a) Xác định vị trí ban đầu của xe máy.
b) Tính quãng đường xe máy đi được trong 3 giây đầu tiên và trong 5 giây cuối cùng.
c) Tính vận tốc trung bình của xe máy trong khoảng thời gian từ $t = 0$ đến $t = 10s$.
d) Xe máy có đổi chiều chuyển động không? Nếu có, tại thời điểm nào?

\loigiai{
a) Vị trí ban đầu của xe máy:
Từ đồ thị, tại $t = 0s$, độ dịch chuyển $d = 20 \text{ m}$. Vậy vị trí ban đầu của xe máy là $20 \text{ m}$.

b) Quãng đường xe máy đi được:
- Trong 3 giây đầu tiên (từ $t=0$ đến $t=3s$):
  Vận tốc trong giai đoạn này là $v_1 = \frac{d_1 - d_0}{t_1 - t_0} = \frac{50 - 20}{3 - 0} = \frac{30}{3} = 10 \text{ m/s}$.
  Quãng đường đi được là $S_1 = |v_1| \times \Delta t_1 = |10| \times 3 = 30 \text{ m}$.
- Trong 5 giây cuối cùng (từ $t=5$ đến $t=10s$):
  Vận tốc trong giai đoạn này là $v_2 = \frac{d_2 - d_1}{t_2 - t_1} = \frac{0 - 50}{10 - 5} = \frac{-50}{5} = -10 \text{ m/s}$.
  Quãng đường đi được là $S_2 = |v_2| \times \Delta t_2 = |-10| \times 5 = 50 \text{ m}$.

c) Vận tốc trung bình của xe máy trong khoảng thời gian từ $t = 0$ đến $t = 10s$:
Độ dịch chuyển tổng cộng là $\Delta d = d_{cuoi} - d_{dau} = 0 - 20 = -20 \text{ m}$.
Thời gian tổng cộng là $\Delta t = 10 - 0 = 10 \text{ s}$.
Vận tốc trung bình là $v_{tb} = \frac{\Delta d}{\Delta t} = \frac{-20}{10} = -2 \text{ m/s}$.

d) Xe máy có đổi chiều chuyển động không?
Vận tốc trong giai đoạn từ $t=0$ đến $t=5s$ là $v_A = \frac{50 - 20}{5 - 0} = \frac{30}{5} = 6 \text{ m/s}$ (chiều dương).
Vận tốc trong giai đoạn từ $t=5$ đến $t=10s$ là $v_B = \frac{0 - 50}{10 - 5} = \frac{-50}{5} = -10 \text{ m/s}$ (chiều âm).
Vì vận tốc đổi dấu từ dương sang âm tại $t = 5s$, nên xe máy có đổi chiều chuyển động tại thời điểm $t = 5s$.
}
\end{ex}

\dangbt{Vận dụng đồ thị độ dịch chuyển -- thời gian của một vật}
\begin{ex}
Một vật chuyển động thẳng có đồ thị độ dịch chuyển – thời gian như hình vẽ.
\begin{center}
\includegraphics[width=0.7\textwidth]{graph_displacement_time_3.png}
\end{center}
a) Xác định thời điểm vật đi qua vị trí có độ dịch chuyển bằng 0 lần thứ hai.
b) Tính vận tốc của vật trong khoảng thời gian từ $t = 0$ đến $t = 2s$ và từ $t = 4s$ đến $t = 6s$.
c) Tính quãng đường vật đi được trong 6 giây đầu tiên.
d) So sánh độ lớn vận tốc của vật trong giai đoạn từ $t = 0$ đến $t = 2s$ và từ $t = 2s$ đến $t = 4s$.

\loigiai{
a) Thời điểm vật đi qua vị trí có độ dịch chuyển bằng 0 lần thứ hai:
Từ đồ thị, vật đi qua vị trí có độ dịch chuyển bằng 0 lần thứ nhất tại $t = 0s$.
Lần thứ hai vật đi qua vị trí có độ dịch chuyển bằng 0 là tại $t = 6s$.

b) Tính vận tốc của vật:
- Trong khoảng thời gian từ $t = 0$ đến $t = 2s$:
  $v_1 = \frac{d_1 - d_0}{t_1 - t_0} = \frac{4 - 0}{2 - 0} = 2 \text{ m/s}$.
- Trong khoảng thời gian từ $t = 4s$ đến $t = 6s$:
  $v_2 = \frac{0 - 4}{6 - 4} = \frac{-4}{2} = -2 \text{ m/s}$.

c) Quãng đường vật đi được trong 6 giây đầu tiên:
- Giai đoạn 1 (từ 0 đến 2s): Độ dịch chuyển $\Delta d_1 = 4 - 0 = 4 \text{ m}$. Quãng đường $S_1 = |\Delta d_1| = 4 \text{ m}$.
- Giai đoạn 2 (từ 2 đến 4s): Độ dịch chuyển $\Delta d_2 = 4 - 4 = 0 \text{ m}$. Quãng đường $S_2 = |\Delta d_2| = 0 \text{ m}$.
- Giai đoạn 3 (từ 4 đến 6s): Độ dịch chuyển $\Delta d_3 = 0 - 4 = -4 \text{ m}$. Quãng đường $S_3 = |\Delta d_3| = 4 \text{ m}$.
Tổng quãng đường vật đi được là $S = S_1 + S_2 + S_3 = 4 + 0 + 4 = 8 \text{ m}$.

d) So sánh độ lớn vận tốc:
- Vận tốc trong giai đoạn từ $t = 0$ đến $t = 2s$: $v_A = 2 \text{ m/s}$. Độ lớn $|v_A| = 2 \text{ m/s}$.
- Vận tốc trong giai đoạn từ $t = 2s$ đến $t = 4s$: $v_B = \frac{4 - 4}{4 - 2} = 0 \text{ m/s}$. Độ lớn $|v_B| = 0 \text{ m/s}$.
Vậy độ lớn vận tốc của vật trong giai đoạn từ $t = 0$ đến $t = 2s$ lớn hơn độ lớn vận tốc trong giai đoạn từ $t = 2s$ đến $t = 4s$. ($2 \text{ m/s} > 0 \text{ m/s}$)
}
\end{ex}

\dangbt{Vận dụng đồ thị độ dịch chuyển -- thời gian của một vật}
\begin{ex}
Một người đi bộ chuyển động thẳng có đồ thị độ dịch chuyển – thời gian như hình vẽ.
\begin{center}
\includegraphics[width=0.7\textwidth]{graph_displacement_time_4.png}
\end{center}
a) Xác định thời điểm người đi bộ bắt đầu chuyển động và vị trí ban đầu của người đó.
b) Tính vận tốc của người đi bộ trong khoảng thời gian từ $t = 0$ đến $t = 3s$ và từ $t = 3s$ đến $t = 6s$.
c) Tính độ dịch chuyển tổng cộng của người đi bộ sau 8 giây.
d) Trong khoảng thời gian nào người đi bộ đứng yên?

\loigiai{
a) Thời điểm người đi bộ bắt đầu chuyển động và vị trí ban đầu:
Người đi bộ bắt đầu chuyển động tại $t = 0s$.
Vị trí ban đầu của người đó là $d = 0 \text{ m}$.

b) Tính vận tốc của người đi bộ:
- Trong khoảng thời gian từ $t = 0$ đến $t = 3s$:
  $v_1 = \frac{d_1 - d_0}{t_1 - t_0} = \frac{6 - 0}{3 - 0} = 2 \text{ m/s}$.
- Trong khoảng thời gian từ $t = 3s$ đến $t = 6s$:
  $v_2 = \frac{6 - 6}{6 - 3} = 0 \text{ m/s}$.

c) Độ dịch chuyển tổng cộng của người đi bộ sau 8 giây:
Tại $t = 8s$, độ dịch chuyển của người đi bộ là $d = 2 \text{ m}$.
Độ dịch chuyển tổng cộng là $\Delta d = d_{cuoi} - d_{dau} = 2 - 0 = 2 \text{ m}$.

d) Trong khoảng thời gian nào người đi bộ đứng yên?
Người đi bộ đứng yên khi đồ thị độ dịch chuyển – thời gian là một đường thẳng nằm ngang (độ dịch chuyển không đổi).
Từ đồ thị, người đi bộ đứng yên trong khoảng thời gian từ $t = 3s$ đến $t = 6s$.
}
\end{ex}
